%%%%%%%%%%%%%%%%%%%%%%%%%%%%%%%%
% Introduction
%%%%%%%%%%%%%%%%%%%%%%%%%%%%%%%%

\section{Introduction}

Many computational systems rely on \emph{approximate set representations}---data structures whose membership queries may err, returning false positives or false negatives with known rates.
Bloom filters~\cite{bf}, perfect hash filters~\cite{phf}, and related probabilistic data structures achieve dramatic space savings by tolerating such errors, and individual structures are well understood.
However, practical systems rarely use a single approximate set in isolation.
Encrypted search indexes, database query planners, and distributed systems routinely compose approximate sets through union, intersection, complement, and difference, and the resulting error behavior has traditionally been analyzed on a case-by-case basis.

\paragraph{Scope and prior work.}
The per-element error model used here is classical.
A membership indicator observed through independent false positive and false negative errors is an asymmetric binary channel~\cite{shannonBSC}, the misclassification model for $2 \times 2$ tables~\cite{bross1954misclassification}, and Warner's randomized response~\cite{warner1965} applied to the indicator of the objective set.
Probabilistic sets, described by per-element membership probabilities, were introduced by Hirota~\cite{hirota1981probabilistic}, and filters with both false positives and false negatives include retouched Bloom filters~\cite{donnet2006retouched} and generalized Bloom filters~\cite{laufer2011generalized}.
This paper claims neither a new model nor new theorems; it collects these ideas into one axiomatic presentation on which the companion papers build.

\paragraph{Contents.}
We state the \emph{Bernoulli set model} for \emph{random approximate sets} as two axioms: element-wise independence of errors and conditional independence of block error rates.
From these axioms we derive the following standard consequences:
\begin{enumerate}[nosep]
    \item the binomial distributions of error counts (false positives, false negatives, true positives, true negatives) and their asymptotic normal limits,
    \item the $k$-fold composition of approximate identity functions, which is a cascade of binary channels whose rates are given by the product of channel transition matrices~\cite{coverThomas}, and
    \item probability intervals for the realized false positive and true positive rates.
\end{enumerate}
The framework is formulated as an abstract data type: any implementation whose membership queries satisfy the Bernoulli axioms inherits these results.
Two qualifications apply.
A Bloom filter satisfies the independence axiom only conditionally on its realized bit array (remark ``\nameref{rem:implementations}'' in \cref{sec:bernoulli_model}), and set-theoretic compositions of Bernoulli sets are higher-order models rather than second-order Bernoulli sets (\cref{sec:higher_order_model}).

\paragraph{Related work.}
The original Bloom filter~\cite{bf} introduced approximate set membership with one-sided error.
Subsequent work has produced a rich family of probabilistic data structures---including cuckoo filters~\cite{cuckooFilter}, quotient filters~\cite{quotientFilter}, xor filters~\cite{xorFilter}, and ribbon filters~\cite{ribbonFilter}---surveyed by Broder and Mitzenmacher~\cite{broderMitzenmacher}, with improved space efficiency or additional functionality.
Bose et al.~\cite{boseBloom} give tight bounds on the Bloom filter false positive rate, Kirsch and Mitzenmacher~\cite{kirschMitzenmacher} show that two hash functions suffice, and Carter and Wegman~\cite{carterWegman} provide the universal hash families underlying many implementations.
Related probabilistic summaries such as the count-min sketch~\cite{countMinSketch} trade exact answers for space efficiency in the streaming setting.
Set operations on Bloom filters (bitwise OR for union, bitwise AND for intersection) are also described in that survey~\cite{broderMitzenmacher}.
General probabilistic methods for randomized algorithms~\cite{mitzenmacherUpfal} provide the analytical toolkit (concentration inequalities, entropy bounds) used throughout.
The present framework is an organizing layer over these tools rather than a replacement for them; the interval arithmetic extension for uncertain rate parameters is developed in~\cite{bernoulliComposition}.

\paragraph{Companion papers.}
The set-theoretic composition of Bernoulli sets---closed-form error rates for complement, union, intersection, and set difference, monoidal structure, relational predicates, and interval arithmetic---is developed in a companion paper~\cite{bernoulliComposition}.
The probability distributions of binary classification measures (precision, recall, accuracy, etc.) induced by the Bernoulli model are derived in~\cite{bernoulliMeasures}.
The joint entropy of error counts and the information-theoretic space complexity of approximate sets are developed in~\cite{bernoulliEntropy}.

\paragraph{Organization.}
\Cref{sec:setalgebra} reviews the algebra of sets.
\Cref{sec:bernoulli_model} introduces the Bernoulli set model and its axioms, including the higher-order composition theorem.
\Cref{sec:characteristics} derives the binomial distributions of error counts and their asymptotic normal limits.
